\section{Formalization architecture and boundary}\label{sec:formalization}

The \texttt{BoundedArityDecoders} Lean~4 package imports the repository's existing libraries
without modifying any upstream library, namely \leanname{Transpose} (\cite{FirstPaper2026}),
\leanname{BoundedStackDepth} (\cite{BoundedStack2026}), \leanname{TapeStorageRates}
(\cite{HighRate2026}), \leanname{SwapRigidity} (\cite{SeventeenthPaper2026}),
\leanname{AdaptiveTransfer} (\cite{AdaptiveTransfer2026}), and \leanname{SharedDecoders}
(\cite{SharedDecoders2026}), all pinned to Lean and Mathlib version \texttt{4.34.1}
\cite{Mathlib2026}.

The formalization is organized into twenty modules across three layers:
\begin{enumerate}[label=(\arabic*)]
\item \emph{Seven formal specification modules}
  (\leanfile{BoundedAritySetCapacity.lean},
  \leanfile{AffineCubeComposition.lean},
  \leanfile{SumsetRankDeficit.lean},
  \leanfile{OperatorDuality.lean},
  \leanfile{SwapCodeAnswers.lean},
  \leanfile{NativeSweepInterface.lean}, and
  \leanfile{DiagnosisAndBarriers.lean}), defining $51$ named `: Prop` specifications covering every
  theorem, lemma, finite control, prior-spine benchmark, and open target of the paper;
\item \emph{Seven primary proof modules} (\texttt{*Proofs.lean}), establishing complete
  kernel-checked proofs of every proved specification; and
\item \emph{Six consequence and representation-bridge modules}
  (\leanfile{SetCapacityConsequences.lean},
  \leanfile{AffineCubeConsequences.lean},
  \leanfile{SumsetRankConsequences.lean},
  \leanfile{DualityAndSwapConsequences.lean},
  \leanfile{MachineSweepConsequences.lean}, and
  \leanfile{DiagnosisAndBridgesConsequences.lean}), proving the unconditional empty/nonempty set
  extensions, explicit constants, all-dimension asymptotics, Mathlib \leanname{TensorProduct}
  intertwining equivalences, discrete Shannon mutual information identities, native
  \leanname{Transpose.TM} symbolic tree extractors, dummy-tape padding, and conditional closures of
  $(\T)$.
\end{enumerate}

Unlike manuscripts with written-only sections, \emph{every} numbered theorem, lemma, proposition,
corollary, and example in this paper is completely proved in Lean~4, including the explicit
$3$-work-tape and $2$-work-tape \leanname{Transpose.TM} transition-table constructions for Stone's
shuffle, Stone's unshuffle, and row shear (\Cref{prop:stone-tm}), the native symbolic causal-cone
tree evaluators (\Cref{lem:tm-cut-cone,thm:sweep-cut-absorbed,thm:sweep-cut-refined}), the
explicit primitive $12$-th root of unity $\zeta_{12} = \exp(2\pi i / 12)$ in $\C$
(\Cref{ex:d12-band}), and the finite $d = 4, 7, 12, 13$ sharpness and characteristic-split
witnesses. Open questions ($(\T)$, the $M=1, \Delta=O(1)$ half-barrier, the macroscopic two-stage
rank target, the $d=5$ three-term threshold, and the periodic Walsh multi-scale target) are defined
as explicit `: Prop` declarations so that all logical equivalences and conditional implications
among them are themselves machine-checked theorems. Where \Cref{prop:diag-high-rate}(a) formalizes
the separated storage-code amplification benchmark from \cite{HighRate2026}, its theorem statement
retains the explicit odd-girth seed hypothesis \lean{HasOddGirthLinearSeed} (matching
\leanname{TapeStorageRates.OddGirthLinearSeed} of \cite[Definition~5.1]{HighRate2026}).
